\documentclass{article}
\usepackage{graphicx} 
\usepackage{amsmath, amssymb}

\title{Travail pratique de mathématiques algèbre}
\author{NTAYINGI ASEKIN'KALI Marie Louise }
\date{Décembre 2024}

\begin{document}
\maketitle
\pagestyle{empty}
\underline{Question}

sous quelles conditions les propriétés suivantes:
\[
(3) \forall x \in A*2, \forall (n,m) \in \mathbb{N}^2, x^n \cdot x^m = x^{n+m}\]

\[
(4) \forall x \in A*2, \forall (n,m) \in \mathbb{N}^2, (x^n)^m= x^{n \times m}\]

peuvent s'écrire comme suit dans le contexte d'un anneau:\\
\[
(3') \forall x \in A*2, \forall (n,m) \in \mathbb{Z}^2, x^n \cdot x^m = x^{n+m}\]

\[
(4') \forall x \in A*2, \forall (n,m) \in \mathbb{Z}^2, (x^n)^m = x^{nm}\]

\underline{\textbf{\textit{Résolution}}}\\ 

les propriétés (3') et (4') peuvent s'écrire dans le contexte d'un anneau avec comme support \(Z^2\) si:\\
Sachant que \( \mathbb{Z} = \{ \ldots, -2, -1, 0, 1, 2, \ldots \} \), il faudrait traiter les exposants négatifs.\\
C'est-à-dire, pour que \( x^n \) soit défini pour \( n < 0 \),\\
on doit définir \( x^{-n} = \frac{1}{x^n} \) lorsque \( x \neq 0 \).\\
Cela nécessite que chaque élément non nul de \( A \) ait un inverse multiplicatif.\\

(e) Il doit exister un inverse multiplicatif \( y \) pour la seconde loi tel que \( x \cdot y = e \), où \( e \) est l'élément neutre.\\

(1) \(\forall x \in A, x \neq 0, \forall k \in \mathbb{Z}, k < 0: x^k = \frac{1}{x^{-k}} \).
(2) \(\forall x \in A, x \neq 0, \exists y \in A, y = x^{-1} : x \cdot x^{-1} \leftrightarrow x \cdot y = e \leftrightarrow x \cdot \frac{1}{x} = e.\)\\

4. Pour \(n < 0\), on doit pouvoir écrire :\\
\((x^n)^m = (x^n)^{-m} \leftrightarrow x^{-n \cdot m}.\)\\
Ce qui signifie que :\\
\(\forall x \in A, x \neq 0, \forall n, m \in \mathbb{Z} : (x^n)^m = x^{n \cdot m}.\)\\

Si \(n = a\) et \(m = b\), alors :\\
\((x^n)^m = (x^a)^b = x^{a \cdot b}.\)\\

Si \(n = a\) et \(m = -b\), alors :\\
\((x^n)^m = (x^a)^{-b} = x^{a \cdot (-b)} = x^{-(a \cdot b)}.\)

2. Il doit exister un élément neutre pour la seconde loi dans le cas des exposants positifs et négatifs, car un anneau \((A, +, \cdot)\) est d'abord un groupe. \((A, +)\), la première loi, admet déjà un élément neutre sur \(A\).\\

(1) \(\exists e \in A, \forall x \in A, x \neq 0 : x \cdot e = x \text{ et } e \cdot x = x.\)\\
(1') \(\exists e \in A, \forall x \in A, x \neq 0 : x^{-1} \cdot e = x^{-1} \text{ et } e \cdot x^{-1} = x^{-1}.\)\\

(2) \(\forall x \in A, x \neq 0.\)\\
(2') \(\exists e \in A, \forall x \in A, x \neq 0 : x^0 = e.\)\\

(1) et (1') :\\
\(\exists e \in A, \forall x \in A, x \neq 0 : (x \cdot e = x \text{ et } e \cdot x = x) \text{ et } (x^{-1} \cdot e = x^{-1} \text{ et } e \cdot x^{-1} = x^{-1}).\)\\

3. Pour \(n < 0\), on doit pouvoir écrire \(x^n\) comme \(x^{-n}\), pour \(m = -n > 0\), ce qui implique que :\\
\(x^n \cdot x^m = x^{n + m}\) doit être vrai pour \(n, m \in \mathbb{Z}.\)
(1) \(\forall x \in A, x \neq 0, \forall n, m \in \mathbb{Z}, m = -n > 0 : x^n \cdot x^m = x^{n+m} = x^{n-n} = x^0 = e \in A.\)\\
(2) \(\forall x \in A, x \neq 0, \forall n, m \in \mathbb{Z} : x^n \cdot x^m = x^{n+m}.\)\\

Cela signifie que si \(n = a\) et \(m = b\) avec \(a, b \in \mathbb{Z}\) :\\

(1) \(x^a \cdot x^b = x^{a+b} = \frac{1}{x^{-a}} \cdot \frac{1}{x^{-b}} = \frac{1}{x^{-(a+b)}} = x^{-(-(a+b))} = x^{a+b}.\)\\
Si \(n = a\) et \(m = -b\) :\\
(2) \(x^a \cdot x^m = x^{a-b} = \frac{1}{x^{-a}} \cdot \frac{1}{x^{-(-b)}} = \frac{1}{x^{-a}} \cdot x^b = \frac{1}{x^{-a+b}} = x^{-(-a+b)} = x^{a-b}.\)

\maketitle
\end{document}
